% Compile with: pdflatex lecture29.tex
% Requires: amsmath, amssymb, amsthm, mathtools, xcolor, mdframed,
%           enumitem, hyperref, pagecolor, microtype

\documentclass[11pt]{article}
\usepackage[margin=1in]{geometry}
\usepackage{amsmath, amssymb, amsthm}
\usepackage{mathtools}
\usepackage{xcolor}
\usepackage{mdframed}
\usepackage{enumitem}
\usepackage{hyperref}
\usepackage{pagecolor}
\usepackage{microtype}

% ---------- Colours ----------
\definecolor{cream}{HTML}{FFFDF5}
\definecolor{defblue}{HTML}{1A4A7A}
\definecolor{thmgreen}{HTML}{1A5C3A}
\definecolor{warnAmber}{HTML}{7A4A00}
\definecolor{dangerRed}{HTML}{7A1A1A}
\definecolor{boxbluebg}{HTML}{EBF3FB}
\definecolor{boxgreenbg}{HTML}{EBF7EF}
\definecolor{boxamberbg}{HTML}{FDF5E6}
\definecolor{boxredbg}{HTML}{FBEBEB}
\pagecolor{cream}

% ---------- Theorem styles ----------
\newtheoremstyle{defstyle}    {8pt}{8pt}{\normalfont}{0pt}{\bfseries\color{defblue}}  {.}{0.5em}{}
\newtheoremstyle{thmstyle}    {8pt}{8pt}{\normalfont}{0pt}{\bfseries\color{thmgreen}} {.}{0.5em}{}
\newtheoremstyle{notstyle}    {8pt}{8pt}{\normalfont}{0pt}{\bfseries\color{warnAmber}}{.}{0.5em}{}
\newtheoremstyle{cautionstyle}{8pt}{8pt}{\normalfont}{0pt}{\bfseries\color{dangerRed}}{.}{0.5em}{}

\theoremstyle{defstyle}
\newtheorem{definition}{Definition}[section]

\theoremstyle{thmstyle}
\newtheorem{theorem}{Theorem}[section]
\newtheorem{corollary}[theorem]{Corollary}
\newtheorem{lemma}[theorem]{Lemma}
\newtheorem{proposition}[theorem]{Proposition}

\theoremstyle{notstyle}
\newtheorem*{remark}{Remark}
\newtheorem*{example}{Example}

\theoremstyle{cautionstyle}
\newtheorem*{caution}{Caution}

% ---------- Boxes ----------
\newmdenv[
  backgroundcolor=boxbluebg, linecolor=defblue, linewidth=2pt,
  topline=false, rightline=false, bottomline=false,
  innerleftmargin=10pt, innerrightmargin=10pt,
  innertopmargin=6pt, innerbottommargin=6pt,
  skipabove=8pt, skipbelow=8pt
]{defbox}

\newmdenv[
  backgroundcolor=boxgreenbg, linecolor=thmgreen, linewidth=2pt,
  topline=false, rightline=false, bottomline=false,
  innerleftmargin=10pt, innerrightmargin=10pt,
  innertopmargin=6pt, innerbottommargin=6pt,
  skipabove=8pt, skipbelow=8pt
]{thmbox}

\newmdenv[
  backgroundcolor=boxamberbg, linecolor=warnAmber, linewidth=2pt,
  topline=false, rightline=false, bottomline=false,
  innerleftmargin=10pt, innerrightmargin=10pt,
  innertopmargin=6pt, innerbottommargin=6pt,
  skipabove=8pt, skipbelow=8pt
]{notebox}

\newmdenv[
  backgroundcolor=boxredbg, linecolor=dangerRed, linewidth=2pt,
  topline=false, rightline=false, bottomline=false,
  innerleftmargin=10pt, innerrightmargin=10pt,
  innertopmargin=6pt, innerbottommargin=6pt,
  skipabove=8pt, skipbelow=8pt
]{cautionbox}

% ---------- Macros ----------
\newcommand{\Z}{\mathbb{Z}}
\newcommand{\N}{\mathbb{N}}
\newcommand{\R}{\mathbb{R}}
\newcommand{\card}[1]{\lvert #1 \rvert}
\newcommand{\Npos}{\mathbb{N}^+}
\newcommand{\fib}[1]{f^{-1}(\{#1\})}

\begin{document}

\begin{center}
  {\Large\bfseries Lecture 29 --- Strict Cardinality and Pigeonhole}\\[4pt]
  {\normalsize CS 173: Discrete Structures \quad$\cdot$\quad $\card{A} < \card{B}$, $\card{A} > \card{B}$}\\[2pt]
  {\small\color{gray} University of Illinois at Urbana-Champaign}
\end{center}
\vspace{1em}\hrule\vspace{1.5em}

% ============================================================
\section{Strict Cardinality}
% ============================================================

\begin{defbox}
\begin{definition}[$\card{A} < \card{B}$]
For any sets $A, B$:
\[
  \card{A} < \card{B}
  \;\iff\;
  \exists f\,(f : A \to B \,\wedge\, f \text{ injective})
  \;\wedge\;
  \forall g\,(g : A \to B \Rightarrow g \text{ not surjective}).
\]
\end{definition}
\end{defbox}

\begin{defbox}
\begin{definition}[$\card{A} > \card{B}$]
For any sets $A, B$:
\[
  \card{A} > \card{B}
  \;\iff\;
  \exists g\,(g : A \to B \,\wedge\, g \text{ surjective})
  \;\wedge\;
  \forall f\,(f : A \to B \Rightarrow f \text{ not injective}).
\]
\end{definition}
\end{defbox}

\begin{notebox}
\begin{remark}
Each definition pairs an existential witness with a universal failure: $A$ fits
inside $B$ but cannot cover it, or covers $B$ but cannot embed in it. The
universal clause is essential --- a single non-surjective function (e.g.\ a
constant) proves nothing on its own.
\end{remark}
\end{notebox}

% ============================================================
\section{Pigeonhole Theorem}
% ============================================================

\begin{thmbox}
\begin{theorem}[Pigeonhole]
For any sets $A, B$:
\begin{enumerate}[label=\textnormal{(\roman*)}, leftmargin=2em, itemsep=2pt]
  \item $\card{A} > \card{B} \;\Rightarrow\; \forall f\,(f : A \to B \Rightarrow f \text{ not injective})$.
  \item $\card{A} < \card{B} \;\Rightarrow\; \forall f\,(f : A \to B \Rightarrow f \text{ not surjective})$.
\end{enumerate}
\end{theorem}
\end{thmbox}

\begin{notebox}
\begin{remark}
Pigeonhole upgrades the universal clause of each definition into a theorem you
can apply once strict inequality is established.
\end{remark}
\end{notebox}

% ============================================================
\section{Quantitative Pigeonhole}
% ============================================================

The basic theorem says some fibre is non-singleton. The quantitative form pins
down \emph{how} crowded that fibre must be.

\begin{thmbox}
\begin{theorem}[Quantitative Pigeonhole]
Let $\card{A} = n \in \Npos$ and $\card{B} = k \in \Npos$. Then for any
function $f : A \to B$:
\[
  \exists b \in B \;\Bigl(\;\card{\{\,a \in A \mid f(a) = b\,\}} \;\geq\; \left\lceil \tfrac{n}{k} \right\rceil\;\Bigr).
\]
\end{theorem}
\end{thmbox}

\begin{notebox}
\begin{remark}
When $k \nmid n$, the bound rewrites as $\left\lceil \tfrac{n}{k} \right\rceil = \left\lfloor \tfrac{n}{k} \right\rfloor + 1$
--- the form written on the board. When $k \mid n$, ceiling and floor agree
and equal $n/k$.
\end{remark}
\end{notebox}

\begin{notebox}
\begin{remark}[Sanity check]
If every fibre had size $< \lceil n/k \rceil$, i.e.\ at most
$\lceil n/k \rceil - 1$, the total $\sum_{b \in B} \card{\fib{b}}$ would be at
most $k\bigl(\lceil n/k \rceil - 1\bigr) < n$ --- but this sum equals
$\card{A} = n$. Contradiction.
\end{remark}
\end{notebox}

% ============================================================
\section{Finite Sets: One Witness Suffices}
% ============================================================

\begin{thmbox}
\begin{proposition}[Finite shortcut]
Let $A, B$ be finite. Then:
\begin{enumerate}[label=\textnormal{(\roman*)}, leftmargin=2em, itemsep=2pt]
  \item If $f : A \to B$ is injective and not surjective, then $\card{A} < \card{B}$.
  \item If $g : A \to B$ is surjective and not injective, then $\card{A} > \card{B}$.
\end{enumerate}
\end{proposition}
\end{thmbox}

\begin{proof}
\textbf{(i)} $f$ witnesses the existential clause. For the universal: if some
$g : A \to B$ were surjective, then $\card{A} \leq \card{B}$ (from $f$) and
$\card{A} \geq \card{B}$ (from $g$) give $\card{A} = \card{B}$, forcing $f$ to
be a bijection --- contradicting non-surjectivity.

\smallskip
\textbf{(ii)} Symmetric: a rival injection $f$ would force $g$ to be a
bijection, contradicting non-injectivity.
\end{proof}

\begin{cautionbox}
\begin{caution}[Infinite sets break the shortcut]
Take $A = B = \N$, $f(n) = 2n$: injective, not surjective, yet
$\card{\N} = \card{\N}$. The identity $g(n) = n$ kills the universal clause.
For $\card{\N} < \card{\R}$ and similar infinite strict comparisons, the
universal clause must be proved directly --- the job of Cantor's diagonal.
\end{caution}
\end{cautionbox}

% ============================================================
\section{Summary}
% ============================================================

\begin{enumerate}[leftmargin=2em, itemsep=2pt]
  \item Strict cardinality $=$ existential witness $\wedge$ universal failure
        of the opposite kind of function.
  \item Pigeonhole supplies the universal clause once strict inequality holds.
  \item Quantitative form: with $\card{A}=n$, $\card{B}=k$, some fibre has size
        at least $\lceil n/k \rceil$.
  \item Finite case: a non-surjective injection witnesses $\card{A} < \card{B}$;
        a non-injective surjection witnesses $\card{A} > \card{B}$.
  \item Infinite case: the universal clause needs a direct proof.
\end{enumerate}

\end{document}